%---------------------------------------------------------------------------%
%->> 开题报告正文，由 origin.md 转换生成
%---------------------------------------------------------------------------%
\chapter{研究背景与研究意义}

\section{研究背景}

人工智能模型中的高性能计算主要由张量算子完成。通用矩阵乘法、卷积和归约等基础算子通常由厂商算子库提供，但现代模型中还包含大量由算子融合、动态形状、稀疏模式和新型数值格式形成的定制算子。对于这些算子，直接编写CUDA或NPU底层程序需要处理线程组织、共享存储、同步、访存合并和流水调度，开发成本较高。

Triton采用以张量块为核心的编程模型，开发者主要描述程序实例处理的数据块及其运算，编译器负责将块级程序映射到具体硬件，从而减少显式线程级编程工作\citep{tillet2019triton}。这一抽象已经成为深度学习框架生成GPU算子的重要技术路径。然而，Triton仅降低了底层编程门槛，并未消除算子优化中的搜索问题。程序仍需确定块大小、warp数量、流水级数、数据布局、归约策略、分块策略等；同一个语义正确的程序，在不同配置下可能产生明显不同的性能。

近年来，大语言模型在代码生成和程序优化方面表现出较强的Triton算子生成能力，能够产出精度正确的Triton算子候选生成能力，为自动生成Triton算子提供了新的可能。但KernelBench的实验表明，即使是先进模型，生成的Triton算子在同时满足正确性和性能要求时仍与手写算子存在较大差距，其评测中的模型能够达到或超过PyTorch基线的任务比例仍不足20\%；通过编译、执行和性能分析反馈进行多轮迭代能够改善结果，说明单次代码生成并不足以解决高性能内核开发问题\citep{ouyang2025kernelbench}。TritonBench进一步构建了包含184个真实Triton算子的评测集，评测集结果表明Triton生成需要同时处理算子语义、编程接口和硬件性能等多重约束\citep{li2025tritonbench}。

另一方面，数据流NPU通常包含处理单元阵列、多级片上存储、DMA引擎和片上互连，其性能依赖计算在空间维度和时间维度上的映射。传统GPU后端中的线程、warp和共享内存抽象不能完整表达NPU上的数据驻留、广播、流水传输和计算通信重叠。因此，在面向NPU的Triton生成过程中，仅搜索源代码层面的块大小和流水级数是不够的，还需要利用NPU后端中间表示、代价模型和诊断信息指导智能体生成与调优。

\section{研究意义}

本文具有以下三方面意义。

第一，在智能程序生成方面，将大语言模型的程序理解能力与编译器的确定性分析能力结合，探索从算子语义描述到高性能Triton程序的自动生成方法，有望降低长尾算子开发与迁移成本。

第二，在智能体与编译器协同方面，将NPU后端的中间表示、合法性约束和性能诊断转化为智能体可理解、可行动的反馈，有望突破仅依赖源代码提示和黑盒运行时间的Triton生成方式。

第三，在协同优化方面，突破传统“智能体只生成代码、编译器只执行优化”的单向关系，使编译器为智能体提供可解释反馈，智能体又能够将优化经验反哺至编译器，实现编译系统从静态规则驱动向“规则、搜索与经验共同驱动”演进。


\chapter{国内外研究现状}

\section{张量编译器与Triton编译技术}

TVM构建了从计算图到张量程序和目标代码的端到端编译体系\citep{chen2018tvm}；Ansor通过自动构造张量程序搜索空间和学习代价模型，降低了人工调度模板依赖\citep{zheng2020ansor}；MLIR通过多层次方言和渐进式Lowering，为异构编译器提供了可扩展基础设施\citep{lattner2021mlir}。在此基础上，TensorIR通过可调度块和可组合变换表示张量程序\citep{feng2023tensorir}，Hidet提出任务映射编程范式\citep{ding2023hidet}，SparseTIR则将可组合格式与变换扩展到稀疏计算\citep{ye2023sparsetir}。这些工作表明，张量编译器的研究重点已由单一循环优化转向可组合程序变换、自动搜索及多层中间表示。

在深度学习框架侧，PyTorch 2利用TorchDynamo和TorchInductor将动态图捕获、图级优化与Triton代码生成结合起来\citep{ansel2024pytorch2}。近期研究进一步关注Triton布局表达和程序变换的可靠性。ASPLOS 2026的Linear Layouts使用有限域上的线性代数统一描述Triton数据布局及其转换\citep{zhou2026linear}；Trinity利用Tile级等式饱和探索张量程序优化空间\citep{park2026trinity}；FuseFlow研究了流式数据流架构上的稀疏算子融合和数据流生成\citep{lacouture2026fuseflow}。这些工作反映出张量编译器正逐渐将布局、融合、数据流和搜索纳入统一优化框架。

但现有Triton体系仍主要围绕GPU执行模型构建，其调度变量、布局语义和代价模型通常隐含SIMT、warp和共享存储等架构假设。对于具有显式数据搬运和空间映射特征的NPU，智能体若只观察源代码和运行时间，难以判断性能瓶颈来自布局、片上存储、DMA流水还是处理单元映射。

\section{智能体代码生成与编译优化}

现有学习型程序优化工作主要包括源代码优化生成、编译参数选择和高性能内核生成。Learning Performance-Improving Code Edits利用大规模代码修改数据学习能够改善性能的程序变换\citep{shypula2024learning}；LLM Compiler针对LLVM IR和汇编程序进行专门训练，说明领域语料与编译任务微调能够增强模型对中间表示及优化行为的理解\citep{cummins2025llmcompiler}；CompilerGym将编译器观察空间、动作空间和奖励函数封装为可交互环境，为强化学习和搜索算法参与编译优化提供了基础\citep{cummins2022compilergym}。

针对GPU内核生成，KernelBench和TritonBench分别从PyTorch算子替换和真实Triton算子生成角度建立了评测体系\citep{ouyang2025kernelbench,li2025tritonbench}。已有研究共同表明，高性能内核生成不能仅依据自然语言和参考代码完成：模型还需要访问编译错误、运行结果、性能计数器以及硬件资源约束。Compiler Generated Feedback for Large Language Models尝试将编译结果重新反馈给模型\citep{grubisic2024feedback}；ASPLOS 2026的LOOPRAG则结合检索、编译测试和性能反馈改进循环变换\citep{zhi2026looprag}，说明外部工具反馈和历史优化案例是提升生成质量的重要途径。

当前方法仍存在三点不足：其一，多数系统以源代码是否可编译作为主要反馈，缺少张量索引、数据布局和数值误差等结构化语义信息；其二，优化经验往往停留在智能体上下文或样例库中，未转化为编译器内部可复用的规则；其三，现有GPU内核生成工作较少考虑数据流NPU中的架构特征。

\section{数据流NPU编译与映射}

针对空间加速器，MAESTRO使用数据中心化指令描述空间—时间映射并估算延迟、能耗和数据复用\citep{kwon2019maestro}；Timeloop建立了统一的架构与映射表示，用于系统化探索DNN加速器设计空间\citep{parashar2019timeloop}；TENET采用关系中心表示统一描述张量运算、处理单元映射和互连关系\citep{lu2021tenet}；CoSA将空间加速器调度建模为带约束的混合整数优化问题\citep{huang2021cosa}；TileFlow进一步针对融合算子建立树形数据流表示和性能分析方法\citep{zheng2023tileflow}。

近年的HPCA和ISCA研究进一步聚焦通信、稀疏性和体系结构—编译协同。HPCA 2025的SoMa研究DNN加速器DRAM通信调度空间\citep{cai2025soma}，CROSS在稀疏和稠密计算内核之间进行编译器驱动选择\citep{liu2025cross}，LEGO研究面向张量应用的空间加速器生成与优化\citep{lin2025lego}；ISCA 2025的ATiM研究Processing-in-DRAM上的张量程序自动调优\citep{shin2025atim}，HYTE采用静态—动态混合方式处理稀疏加速器的分块问题\citep{li2025hyte}，NUPEA则通过非均匀处理单元访问优化空间数据流中的关键访存\citep{ghosh2025nupea}。这些工作说明，数据流架构性能越来越依赖编译器对计算映射、存储容量、通信调度和动态不规则性的联合处理。

总体来看，现有研究分别在Triton生成、张量编译和NPU数据流映射方面取得了进展，但尚缺少一个以Triton为统一接口、连接智能体生成、NPU后端IR级反馈和编译器知识演进的完整协同体系。


\chapter{研究目标与研究问题}

本文的总体目标是构建一个面向数据流NPU Triton后端的算子生成与协同优化原型系统。系统接受PyTorch参考实现、算子描述、输入形状和目标硬件配置，由智能体生成Triton候选程序，经NPU Triton编译器后端完成语义验证、IR级优化、映射和执行，再利用编译与运行结果驱动智能体和编译器共同优化。与仅以自然语言提示生成Triton代码的方式不同，本文希望将算子语义、Triton程序结构、NPU后端IR和硬件资源约束纳入统一闭环，使智能体不只是生成可运行代码，而是能够在编译器反馈的约束下持续修复语义错误、调整调度策略并积累可复用优化经验。

本文首先要解决从参考实现到Triton程序的可靠生成问题。对于矩阵乘、归约、归一化、注意力和融合算子等典型负载，系统需要能够识别输入输出张量之间的索引关系、并行维度、归约维度、广播关系、边界条件和数值误差要求，并将这些信息转化为智能体生成Triton代码时可使用的中间规划。该目标强调正确性优先：生成程序不仅要通过编译，还应在随机输入、边界形状和不同精度条件下与PyTorch参考结果保持一致，为后续性能优化提供可信基础。

其次，本文要解决面向数据流NPU的性能调优问题。NPU后端的优化对象并不限于Triton源代码中的块大小、warp数量或流水级数，还包括片上存储容量约束、DMA搬运、数据驻留、计算与通信重叠、处理单元映射以及Pass选择等因素。因此，本文拟将NPU Triton编译器后端暴露的中间表示、合法性检查结果、资源统计和性能诊断转化为智能体可理解的结构化反馈，使智能体能够根据具体瓶颈区分“程序语义错误”“调度配置不合理”“布局或映射低效”和“后端优化机会不足”等不同情况，并采取相应的代码修改、参数搜索或Pass配置策略。

再次，本文要探索智能体与编译器之间的双向协同机制。传统自动调优通常把编译器视为黑盒执行器，智能体生成方法也多将编译错误和运行时间作为一次性反馈。本文希望进一步把多轮生成、编译、执行和修复过程中形成的优化轨迹沉淀为可检索、可验证、可复用的知识：一方面为智能体后续生成提供案例和先验，另一方面为编译器的Pass选择、代价模型校准和局部优化规则提供依据。该目标要求经验沉淀具有边界条件和置信度，不能将偶然适用于单一形状或单一硬件配置的策略直接提升为全局规则。

最终，本文预期形成一条从“算子语义描述”到“高性能NPU Triton实现”再到“编译器知识更新”的闭环技术路线，并通过代表性算子实验验证其有效性。评价重点包括生成代码的编译成功率和功能正确率、相对PyTorch或现有Triton实现的性能收益、达到目标性能所需的候选数量和编译次数、IR反馈对搜索效率的提升，以及跨形状、跨算子复用优化经验的能力。

围绕这一目标，拟回答三个核心研究问题：

\textbf{RQ1：} 如何将算子语义、Triton编程约束和硬件性能知识组织为智能体可使用的结构化信息，使其生成的程序同时满足正确性与性能要求？

这一问题关注生成起点的可靠性。智能体若直接根据自然语言描述或PyTorch代码输出Triton程序，容易遗漏索引边界、广播规则、归约顺序、掩码处理和数值精度等关键约束，也难以显式表达数据复用和分块策略。本文需要研究一种适合张量算子的中间规划表示，用于描述输入输出张量关系、计算维度分类、分块层次、临时数据生命周期、候选融合策略和误差容忍范围，并在此基础上约束Triton代码生成。该问题还包括如何构造检索样例、模板和错误修复案例，使智能体在生成时能够利用已有高质量程序和失败经验，而不是在巨大代码空间中盲目试错。

\textbf{RQ2：} 如何将NPU Triton编译器后端的中间表示、Pass诊断、资源约束和性能模型组织为智能体可使用的反馈，使代码生成、调度配置和IR级优化能够联合调优？

这一问题关注优化闭环的有效性。对于数据流NPU，同一段Triton代码在不同布局、分块、映射和Pass序列下可能产生差异显著的执行结果，仅依赖最终运行时间无法解释性能瓶颈，也难以指导下一步修改。本文需要研究如何从后端IR、Pass日志、资源占用、DMA传输量、流水停顿、片上存储冲突、处理单元利用率和性能模型误差中抽取可行动信息，并将其映射为智能体的修改建议或搜索约束。例如，当瓶颈来自片上存储溢出时，应优先调整分块和数据驻留；当瓶颈来自DMA与计算重叠不足时，应优先改变流水和搬运策略；当瓶颈来自Pass选择不当时，则需要把优化目标从源代码修改扩展到编译流程配置。

\textbf{RQ3：} 如何将单次编译和运行反馈转化为可跨算子复用的知识，并在不破坏编译器正确性和可复现性的前提下反哺Pass选择、代价模型和优化规则？

这一问题关注协同系统的长期演进。单个算子的调优结果如果只保存在一次对话上下文或一次实验日志中，其价值难以迁移到后续任务；但如果未经验证就转化为编译器规则，又可能引入错误优化或性能退化。本文需要研究优化经验的表示、聚合和门控机制，将“算子特征、形状范围、硬件配置、候选程序、Pass序列、性能结果和失败原因”统一记录，进而提炼出可检索案例、调度先验、Pass推荐策略和候选重写规则。同时，需要通过差分测试、变形测试、资源约束检查和必要的等价性验证控制规则污染，保证智能体反哺编译器的过程可解释、可回滚、可复现。


\chapter{主要研究内容及创新设计}

\section{面向NPU后端协同调优的Triton算子生成智能体}

首先构建面向张量算子的工具调用型智能体。其输入包括自然语言算子描述、PyTorch参考实现、输入输出形状、数据类型、误差范围和目标硬件信息，输出为Triton程序、调度配置以及可选的编译Pass配置建议。

为避免智能体直接从参考代码跳转到Triton源代码，本文拟引入中间规划表示\texttt{KernelPlan}。该表示描述：

\begin{enumerate}
    \item 输入输出张量及其索引关系；
    \item 并行维度、归约维度和广播维度；
    \item 分块层次及边界处理方式；
    \item 数据驻留、复用和中间结果生命周期；
    \item 候选融合、向量化和流水策略。
\end{enumerate}

\nocite{jia2019taso}
智能体首先生成并检查\texttt{KernelPlan}，再将其转换为Triton代码，使语义推理与代码表面形式解耦。在候选生成阶段，引入由Triton模板、历史优质程序、编译错误、IR诊断及失败修复案例构成的检索库。针对一个算子生成多个具有不同分块和布局策略的候选，通过编译、差分测试、IR级合法性检查和性能测量进行筛选。

智能体迭代过程划分为两个阶段。第一阶段以正确性为主，利用编译错误、越界信息、随机输入测试和PyTorch参考结果修复程序；第二阶段以性能为主，利用NPU后端暴露的片上存储占用、数据搬运量、DMA流水停顿、处理单元利用率和执行时间调整程序、调度参数及Pass配置。这样可以减少智能体在语义错误尚未解决时盲目搜索性能参数的问题。

本部分的创新点在于：提出一种“算子语义规划—Triton生成—编译验证—IR反馈—性能修复”的分层智能体，而非单次提示式代码生成；同时将NPU后端的IR诊断、失败案例及其根因作为一等知识，支持面向错误类型和性能瓶颈类型的检索与定向修复。


\section{智能体反哺编译器的双向协同机制}

第二部分是本文区别于一般Triton生成方法的核心。拟设计统一的“编译反馈总线”，将编译器输出的原始日志、IR变化和性能分析结果转换为结构化反馈，包括：

\begin{itemize}
    \item 语法与类型错误；
    \item 索引越界、依赖冲突和布局不一致；
    \item 各级片上存储占用和bank冲突；
    \item DMA传输量及计算通信重叠比例；
    \item 处理单元利用率、片上网络压力和流水停顿；
    \item 各编译Pass前后的IR变化及性能预测误差。
\end{itemize}

在\textbf{快速闭环}中，智能体根据这些信息修改当前算子的Triton代码、调度参数或Pass配置，实现实例级优化。

在\textbf{慢速闭环}中，系统聚合多个算子的优化轨迹，分析“程序特征—编译决策—硬件表现”之间的关系。智能体可向编译器提出三类反馈：一是针对特定程序特征的调度先验，例如归约维度较大时优先采用树形归约；二是Pass选择和排序建议；三是新的局部重写或融合规则。

为防止智能体产生的规则破坏编译器正确性，所有候选规则均需经过多级门控：

\begin{enumerate}
    \item 在历史与自动生成程序上执行差分测试；
    \item 使用随机形状和边界输入进行变形测试；
    \item 对适合的局部变换进行等价性或翻译验证；
    \item 检查存储、带宽和依赖约束；
    \item 只有在多类负载上稳定获益的规则才进入全局规则库。
\end{enumerate}

低置信度经验仅保存在单算子或单架构缓存中，高置信度经验才能提升为编译器规则。与此同时，编译器可将失败调度、反例及性能模型误差转化为智能体的训练样例和检索文档，从而形成真正的双向知识流。

该协同机制包含三个层次：实例级反馈解决当前代码问题，Pass级反馈改善优化流程，规则级反馈促进编译器长期演进。它不同于仅搜索块大小的传统自动调优，也不同于只用编译器报错修复代码的智能体方法。


\chapter{拟解决的关键问题}

\section{算子语义到Triton程序的可靠表达问题}

智能体生成Triton算子时，首先面临的问题不是单纯的代码补全，而是如何准确表达算子语义。PyTorch参考实现中隐含的张量维度关系、广播规则、归约范围、边界掩码和数值误差要求，如果不能被显式抽取和检查，生成的Triton程序即使能够编译，也可能在特殊形状、边界输入或低精度场景下产生错误结果。尤其是在归约、归一化、注意力和融合算子中，性能优化往往会改变计算顺序、临时结果存储方式和中间精度，进一步增加语义验证难度。

因此，本文需要解决从参考实现到\texttt{KernelPlan}再到Triton代码的可信转换问题。一方面，\texttt{KernelPlan}应能描述输入输出索引关系、并行维度、归约维度、分块方式、数据驻留和边界处理等信息，使智能体的生成过程受到结构化语义约束；另一方面，系统需要结合编译检查、随机测试、边界形状测试和误差阈值判断，验证候选程序是否真正等价于参考实现。该问题是后续性能调优的前提，只有保证语义正确，IR反馈和运行时间反馈才具有优化意义。

\section{NPU后端反馈到智能体动作的映射问题}

数据流NPU的性能瓶颈通常不只体现在最终运行时间上，还可能来自片上存储容量不足、DMA搬运过多、计算与通信重叠不足、处理单元利用率偏低、片上网络压力过大或某些Pass引入了低效布局。若智能体只能观察编译是否成功和内核运行时间，就难以判断下一步应修改Triton源代码、调整调度参数，还是改变编译Pass配置。

本文需要将NPU Triton编译器后端产生的IR、Pass诊断、资源统计和性能模型结果转化为智能体可理解、可执行的反馈。关键在于建立“诊断信息—性能瓶颈—优化动作”之间的对应关系：例如，将片上存储溢出关联到分块缩小或中间结果重计算，将DMA停顿关联到数据预取和流水调整，将处理单元利用率不足关联到并行维度选择和映射策略修改，将Pass前后IR变化与性能预测误差关联到Pass选择或排序调整。只有将后端反馈结构化，智能体才能从盲目试错转向有依据的定向修复和调优。

\section{代码、调度参数与IR优化的联合搜索问题}

面向NPU后端的Triton优化并不是单一维度的参数搜索。Triton代码形态、块大小、流水级数、数据布局、片上存储分配、后端映射策略和Pass序列之间存在强耦合：源代码层面的一个分块选择可能改变IR结构，IR结构又会影响Pass适用性和硬件映射结果；某个Pass在一个形状上有效，也可能在另一个形状上导致额外搬运或资源冲突。直接枚举所有候选会造成搜索空间膨胀，难以在可接受编译和运行开销内获得高性能实现。

因此，本文需要研究分层联合调优方法。首先由智能体基于算子语义和检索案例生成较高质量的初始候选，避免从完全随机空间开始搜索；其次利用编译器静态分析和IR合法性检查尽早过滤越界、依赖冲突、资源超限和明显低效的候选；最后针对少量可行候选进行局部调参、Pass配置和性能测量。该问题的核心是把智能体的生成能力、编译器的确定性分析能力和运行反馈结合起来，在正确性、性能收益和搜索成本之间取得平衡。

\section{跨算子优化经验的可信复用问题}

本文不仅关注单个算子的生成和调优，还希望将多轮优化过程中形成的经验反哺给智能体和编译器。然而，优化经验具有明显的适用边界：某一归约策略可能只适用于特定维度规模，某一数据驻留策略可能依赖片上存储容量，某一Pass顺序可能只对特定IR形态有效。如果缺少条件标注和验证机制，系统可能把偶然有效的策略错误推广到其他算子或硬件配置中，造成性能退化甚至错误优化。

因此，需要为优化轨迹建立可复用但受约束的知识表示，记录算子类型、形状范围、数据类型、硬件配置、候选程序特征、Pass序列、性能结果和失败原因等信息。对于低置信度经验，可作为智能体检索案例或局部调度先验；对于希望进入编译器规则库的经验，则必须经过差分测试、变形测试、资源约束检查以及必要的等价性验证。该问题直接关系到双向协同机制能否长期有效：智能体反哺编译器不能只是追加经验样例，而应形成可解释、可回滚、可复现的知识更新流程。


\chapter{研究方法与技术路线}

本文采用“人工高性能样例归纳—智能体生成—编译器反馈调优—经验反哺”的研究方法。其基本思路是：首先从面向NPU架构的手写高性能Triton算子中提取可复用的实现模式和优化技巧，再将这些经验组织为智能体可检索、可调用的知识；随后利用算子语义规划约束Triton代码生成，并通过NPU Triton编译器后端的IR、诊断和性能反馈进行多轮修复与调优；最后将调优轨迹沉淀为跨算子经验，反哺智能体生成、Pass选择和编译器优化规则。

\section{面向NPU的高性能Triton算子收集与模式归纳}

第一步收集和整理针对NPU架构编写的高性能Triton算子。样例来源包括已有NPU后端测试用例、工程实践中的手写算子、公开Triton算子实现以及根据典型模型负载补充实现的基准算子。算子类型覆盖矩阵乘、卷积、Softmax、LayerNorm、RMSNorm、归约、逐元素融合、注意力以及数据搬运密集型算子。对每个算子统一记录PyTorch参考实现、Triton源代码、输入输出形状、数据类型、误差要求、目标NPU配置、编译Pass序列、IR快照、资源统计和性能结果。

在此基础上，对手写高性能算子进行结构化分析，抽取不同类型算子的通用实现优化技巧。对于计算密集型算子，重点分析分块规模、数据复用、累加精度、计算阵列映射和中间结果驻留方式；对于归约和归一化算子，重点分析归约树组织、跨块归约、临时缓冲使用、边界掩码和数值稳定性处理；对于融合算子，重点分析算子融合边界、中间结果消除、访存合并和寄存器或片上存储压力；对于注意力类算子，重点分析分块Softmax、在线归一化、KV数据复用和流水搬运策略。

归纳得到的优化技巧不直接作为自然语言提示堆叠给模型，而是整理为三类知识：一是算子级模板，描述某类算子的基本Triton实现骨架；二是策略级规则，描述在特定形状、数据类型和NPU资源约束下优先采用的分块、布局、流水和数据驻留策略；三是失败修复案例，记录常见编译错误、越界访问、资源超限和性能瓶颈及其修复方式。这样可以为后续智能体生成提供可检索的工程经验基础。

\section{基于KernelPlan的算子语义规划与候选生成}

第二步构建从参考实现到Triton候选程序的生成流程。系统首先解析PyTorch参考实现或算子描述，提取输入输出张量、索引关系、并行维度、归约维度、广播关系、边界条件和误差要求，形成中间规划表示\texttt{KernelPlan}。\texttt{KernelPlan}用于分离算子语义和代码表面形式，使智能体在生成Triton程序前先明确“算什么、按什么维度并行、哪些数据需要复用、哪些边界需要掩码处理”。

在候选生成阶段，智能体根据\texttt{KernelPlan}检索前一阶段建立的算子模板、优化技巧和失败修复案例，生成多个Triton候选程序及其调度配置。候选之间应体现不同的优化取向，例如不同分块形状、不同数据驻留策略、不同归约组织方式、不同流水级数和不同Pass配置建议。生成结果不仅包括源代码，还包括候选策略说明，便于后续根据编译器反馈定位问题和选择修改方向。

\section{NPU后端IR反馈驱动的联合调优}

第三步将生成候选接入NPU Triton编译器后端，建立编译、验证、性能测量和反馈修复闭环。每个候选首先经过语法检查、类型检查、边界检查、差分测试和误差验证，筛除不能正确执行的程序；随后导出后端IR、Pass诊断、片上存储占用、DMA搬运量、处理单元利用率、片上网络压力、流水停顿和性能模型预测等信息。

反馈解析模块将上述信息转换为智能体可执行的调优建议。例如，当IR显示中间结果生命周期过长导致片上存储压力过大时，系统可建议改变分块、减少融合范围或采用重计算；当性能统计显示DMA等待时间较高时，系统可建议调整预取、双缓冲或流水级数；当处理单元利用率不足时，系统可建议改变并行维度或映射策略；当Pass前后IR变化导致性能预测下降时，系统可调整Pass选择和排序。通过这种方式，调优过程不再只依赖最终运行时间，而是利用后端结构化信息进行定向优化。

为控制搜索成本，本文采用分层搜索策略。首先由智能体和检索库产生少量高质量初始候选；其次利用编译器合法性检查、资源约束和性能模型过滤明显不合理的配置；最后对保留下来的候选进行局部参数搜索和真实执行测量。该策略将智能体的经验生成能力、编译器的确定性分析能力和运行反馈结合起来，减少无效编译和盲目试错。

\section{协同知识库构建与反哺机制}

第四步构建协同优化知识库，将生成、编译、运行和修复过程中产生的信息统一记录。知识库条目包括算子语义特征、输入形状范围、数据类型、目标NPU配置、Triton代码特征、\texttt{KernelPlan}、调度参数、Pass序列、IR诊断、性能指标、失败原因和最终修复方式。知识库一方面服务于智能体检索增强，使后续相似算子能够复用已有模板和调优经验；另一方面服务于编译器优化，使稳定有效的经验能够转化为Pass选择先验、代价模型校准样本或候选局部重写规则。

为了避免经验污染，本文将知识反哺划分为不同置信度层次。低置信度经验仅作为智能体提示和相似案例，不直接改变编译器行为；中等置信度经验可用于Pass推荐或搜索空间剪枝；高置信度经验需要经过多形状、多算子和多输入测试，并通过资源约束检查和必要的等价性验证后，才能沉淀为编译器规则或默认调度策略。通过这种门控机制，系统既能利用智能体发现的优化机会，又能保持编译器优化过程的正确性和可复现性。

综上，本文的总体技术路线为：

\begin{quote}
NPU手写高性能Triton算子收集\\
→ 算子实现模式与优化技巧抽取\\
→ 模板、策略规则和失败案例知识库构建\\
→ PyTorch参考实现或算子描述解析\\
→ \texttt{KernelPlan}语义规划生成\\
→ 检索增强的Triton候选生成\\
→ NPU Triton后端编译、IR导出与合法性验证\\
→ 资源诊断、性能预测与执行测量\\
→ 智能体修复、调度搜索和Pass配置调整\\
→ 优化轨迹沉淀与编译器反哺
\end{quote}


\chapter{实验方案与评价指标}

本文实验围绕“高性能样例归纳、智能体生成、NPU后端反馈调优、协同知识反哺”四个环节展开，目标是验证所提出方法是否能够提高Triton算子生成正确率、减少调优搜索成本、提升NPU后端执行性能，并形成可跨算子复用的优化经验。

\section{高性能算子样例与优化技巧抽取实验}

第一类实验验证手写高性能Triton算子收集和优化技巧归纳的有效性。实验对象包括矩阵乘、卷积、Softmax、LayerNorm、RMSNorm、归约、逐元素融合、注意力以及数据搬运密集型算子。对于每类算子，整理若干手写高性能实现和朴素Triton实现，比较二者在代码结构、分块策略、数据驻留、流水组织、边界处理、Pass配置和NPU资源使用上的差异。

评价指标包括算子类型覆盖率、模板覆盖率、可归纳优化技巧数量、失败案例覆盖率、不同优化技巧在样例中的出现频次，以及由模板和策略规则重构出的候选程序相对原手写实现的性能接近程度。该实验重点回答：收集到的手写算子是否能够形成稳定的工程经验，归纳出的模板和策略是否足以支撑后续智能体生成，而不是只对个别样例有效。

\section{基于KernelPlan的Triton候选生成实验}

第二类实验验证\texttt{KernelPlan}和检索增强对Triton代码生成正确率的提升。实验选取KernelBench、TritonBench中的代表性任务，并补充面向NPU后端的归一化、归约、融合和注意力算子。每个任务给定PyTorch参考实现、输入形状、数据类型和误差要求，由系统生成\texttt{KernelPlan}、Triton候选程序和调度配置。

对照方法包括直接使用通用代码模型生成Triton程序、仅使用自然语言提示的智能体、无\texttt{KernelPlan}的检索增强智能体、无手写样例知识库的智能体，以及本文完整生成方法。评价指标包括编译成功率、功能正确率、\texttt{pass@k}、\texttt{correct@k}、边界形状通过率、数值误差通过率、生成候选数量、首次正确候选所需轮次和生成开销。通过消融比较，分析\texttt{KernelPlan}对语义正确性的贡献，以及手写高性能样例知识库对候选质量的贡献。

\section{NPU后端IR反馈驱动的联合调优实验}

第三类实验验证NPU后端IR反馈、资源诊断和性能模型对调优效率和最终性能的作用。实验将已通过正确性验证的Triton候选接入NPU Triton编译器后端，记录后端IR、Pass诊断、片上存储占用、DMA搬运量、处理单元利用率、片上网络通信量、流水停顿和性能预测结果，并据此驱动智能体修改代码、调度参数和Pass配置。

对照方法包括仅使用初始候选不调优、仅基于运行时间进行黑盒搜索、仅使用静态规则调参、仅使用源代码层面反馈，以及本文的IR反馈联合调优方法。评价指标包括内核延迟、有效吞吐率、相对PyTorch或朴素Triton基线的加速比、相对手写高性能算子的性能比例、片上存储利用率、片外数据访问量、DMA与计算重叠比例、处理单元利用率、片上网络压力、达到目标性能所需候选数量、编译次数和搜索总时间。若能够获得实际NPU硬件，则以硬件计数器和真实运行时间为主；否则采用NPU后端性能模型或周期级模拟器，并使用关键微基准对模型误差进行校准。

\section{协同知识反哺与泛化实验}

第四类实验验证优化经验能否跨形状、跨算子和跨配置复用，并考察反哺编译器的可靠性。实验将一部分算子和形状作为训练或经验积累集合，将另一部分未见形状、相似算子和不同NPU资源配置作为测试集合，比较无知识库、仅使用检索案例、仅更新Pass推荐、仅更新代价模型和完整双向闭环等配置。

评价指标包括跨形状性能保持率、未见算子的首次正确候选轮次、达到同等性能所需编译次数、规则或策略命中率、Pass选择准确率、代价模型预测误差、经验复用带来的搜索空间缩减比例，以及规则引入后的性能退化率和正确性失败率。对于拟进入编译器规则库的优化经验，还需统计差分测试通过率、变形测试通过率、资源约束检查通过率和可回滚记录完整性。该实验重点验证智能体反哺编译器是否能够稳定带来收益，并避免将偶然有效的经验错误推广为全局规则。

所有性能实验均采用预热、多次重复和中位数统计，并报告方差、异常值处理方式和测试环境配置。对于生成和调优实验，记录完整的提示、检索案例、\texttt{KernelPlan}、候选代码、Pass序列、IR诊断和性能结果，保证实验过程可复现、错误可定位。


\chapter{可行性分析与风险控制}

Triton、LLVM/MLIR以及多个张量编译和加速器建模框架均提供开源实现，可以支撑原型开发。本文不以重新训练基础大模型或构建完整工业级NPU软件栈为目标，而是重点研究工具调用型智能体、NPU后端反馈接入、Pass配置策略和协同反馈机制，研究范围较为可控。

主要风险包括模型生成稳定性不足、NPU后端接口受限和搜索成本过高。对应措施为：使用确定性编译器检查和模板回退保证系统可用；采用后端日志、IR快照、参数化NPU模型及模拟器降低对专有硬件的依赖；利用解析模型剪枝、历史缓存和分层搜索控制调优成本。


\chapter{研究进度安排}

第1—2个月：完成文献调研、需求分析、算子数据集和评测环境建设。

第3—5个月：实现KernelPlan、检索增强、Triton生成和编译反馈修复智能体。

第6—8个月：完成NPU后端IR解析、性能诊断接入、调度参数搜索和Pass配置推荐机制。

第9—10个月：实现反馈总线、经验库、Pass推荐、调度先验沉淀和优化规则验证机制。

第11个月：开展基线比较、消融实验、跨形状及跨算子泛化实验。

第12个月：整理实验数据，完成论文撰写、系统文档和成果总结。


\chapter{预期成果与创新点总结}

预期形成一个支持代表性深度学习算子的智能体—Triton编译器协同优化原型、一套面向NPU后端IR反馈的Triton算子生成与联合调优方法、一套可复现的算子测试集和优化轨迹数据，并完成研究生毕业论文及相关学术论文。

本文的两项主要创新可概括为：

\begin{enumerate}
    \item \textbf{语义规划与IR反馈驱动的Triton算子生成智能体。} 通过KernelPlan分离算子语义推理与代码生成，并融合NPU后端IR诊断、资源约束和性能反馈，联合调整Triton代码、调度参数和Pass配置。
    \item \textbf{智能体反哺编译器的可信双向闭环。} 将跨算子优化经验转化为调度先验、Pass策略和经验证的编译规则，并通过差分测试、等价性验证和资源门控保证可靠性。
\end{enumerate}
